This chapter analyses how Aether operationalises the architectural connection between edge computing and Private 5G, clarifying their structural interdependence. Rather than treating connectivity and computation as separate layers, Aether implements them through a unified cloud-native control plane governing both network and compute resources \cite{schrammen_edge_computing}. This exposes the dependency between deterministic wireless networking and distributed compute that underpins reliable edge architectures \cite{harmatos_5g_tsn}.

Aether treats Private 5G as a software-defined system. SD-Core and SD-RAN are implemented as microservice-based components with declarative lifecycle semantics \cite{onf_aether_whitepaper, aether_sd_core}, mirroring Kubernetes’ operational model in which reconciliation loops enforce a desired state and enable updates without service disruption. By running SD-Core within a Kubernetes-managed environment, Aether collapses the divide between network operations and application orchestration, subjecting network functions to the same lifecycle and scheduling mechanisms as edge workloads \cite{vachuska_sd_ran, aether_sd_core}. The network therefore becomes part of the same distributed system that governs edge applications.

The relevance to edge computing is most evident in Aether’s compute placement model. Edge architectures require predictable proximity between devices, compute nodes and the network boundary, conditions our prototype \cite{edge_computing_report} achieved only by assuming an uncongested local network and stable links. Aether formalises this by aligning compute placement with 5G topology. Multi-access Edge Computing (MEC) workloads execute near the user plane, typically co-located with DU infrastructure, shortening data paths and placing edge logic within deterministic latency bounds rather than best-effort wireless behaviour \cite{onf_aether_whitepaper}.

This placement model reveals a reciprocal dependency between Private 5G and edge computing. Edge workloads rely on Private 5G for mobility stability, jitter control and QoS isolation, properties that cannot be structurally guaranteed in Wi-Fi or public mobile networks \cite{onf_aether_whitepaper, private_5g_report}. Private 5G, in turn, depends on edge compute to support domain-specific processing close to the radio layer, where filtering, analytics and control loops cannot tolerate unnecessary round-trip delays. MEC and Private 5G therefore form the execution substrate on which edge applications depend, with Aether providing coherent orchestration through Kubernetes primitives such as scheduling, service abstraction and declarative lifecycle management.

Integrating edge computing with Private 5G remains complex. Spectrum licensing, RAN configuration and SD-Core operational overhead introduce constraints that Aether mitigates through automation but cannot eliminate, as radio planning remains inherently physical \cite{private_5g_report}. Its value lies in enforcing architectural constraints through deterministic connectivity, structured compute placement and unified orchestration, enabling edge systems to scale beyond controlled prototype settings. This shifts complexity from ad-hoc deployment decisions into a managed, declarative control plane, reflecting broader trends in system-level reasoning supported by automated and model-assisted design workflows \cite{openai_llm_drafting}.

Aether demonstrates that edge computing and Private 5G are not loosely coupled technologies \cite{harmatos_5g_tsn}. Their relationship is necessarily architectural: distributed compute requires deterministic communication, while Private 5G becomes operationally meaningful only when paired with an edge execution layer \cite{diekmann_ericsson_5g}. In this model, network and compute function as inseparable components of a single distributed system spanning radio, edge and application layers \cite{onf_aether_whitepaper}.
